\section{Open Questions}
Each is grounded both in something Kalidasan et al.\ leave unresolved and
in something this replication surfaced. All five are intended as
research prompts, not paper "future work" boilerplate.

\subsection*{Q1 --- DyP peroxidase: real K279a specificity or ontology artifact?}
\textbf{Basis.} This replication directly contradicts C4 at the gene level:
the DyP gene is present in all four strains under modern RASTtk, but the paper
reported an encapsulin/DyP subsystem exclusive to K279a. The paper itself
never performed a DyP activity assay to confirm K279a-specific function.

\textbf{Next steps.} (1) Re-run classic RAST offline via SEEDtk on the four
assemblies to test whether the 2018 subsystem call reproduces. (2) Search
each genome for the encapsulin shell protein (Pfam PF04454) and ferritin-like
partner (Pfam PF00210) independently of DyP --- the K279a-only signal, if
biological, may live in the shell/partner rather than DyP itself.
(3) Wet-lab: H$_2$O$_2$-dependent ABTS oxidation activity on iron-depleted
lysates of all four strains, gel-filtration into $>200$\,kDa fraction to test
for encapsulation. (4) Cryo-EM / negative-stain EM on K279a to visualise
encapsulin nanocompartments.

\subsection*{Q2 --- What actually is the S.\ maltophilia catechol siderophore?}
\textbf{Basis.} The paper reports catechol-class by Arnow's ("data not shown")
but never identifies the molecule; the paper's own literature citation to
"hydroxamate-type ornibactin" is inconsistent with its catechol result and the
inconsistency is unresolved. This replication also did not run any biosynthetic-
cluster prediction on the four genomes --- a real missed lever.

\textbf{Next steps.} (1) antiSMASH v7/8 across all 4 genomes; enumerate
catecholate NRPS clusters and check for entA/B/E/F homologs. (2) LC-MS/MS on
iron-depleted CS17 supernatant (per paper conditions) after XAD-16 solid-phase
extraction; compare vs enterobactin, 2,3-DHB, and DHB-serine standards.
(3) Delete the top-scoring NRPS cluster in K279a and repeat CAS/Arnow's
--- loss of activity confirms the source cluster. (4) Cross-strain: does
R551-3 have a truncated/absent version of the same cluster, consistent with
its weaker CAS signal?

\subsection*{Q3 --- Clinical $\gg$ environmental: real ecotype, sample-size, or primer artifact?}
\textbf{Basis.} PCR-negative rates of 26.9\%--53.8\% in clinical isolates for
HmuV/FeSS/FCR/Htp, and 0\% amplification of 9-of-17 targets in only 5
environmental isolates, are equally consistent with (a) real gene loss in an
environmental subclade, (b) sample-size noise (n$=$5), or (c) primer failure
(paper Table~3 primers were designed on the 4 in-silico genomes; SNP density
between clinical and environmental \textit{S.\ maltophilia} clades is high).
The paper does not attempt an in-silico primer coverage check; nor did this
replication.

\textbf{Next steps.} (1) Pull all $\gtrsim 200$ public \textit{S.\ maltophilia}
BV-BRC genomes; test the 17-target set for statistically significant
clinical/environmental gene-presence differences (Fisher's exact + BH).
(2) In-silico primer coverage: align each Table~3 primer pair against the
full genome panel; count perfect matches, 3$'$-end mismatches, and predicted
amplicon size. (3) Redesign degenerate primers where coverage $<95\%$ and
reanalyze the 109-isolate panel. (4) Core-genome phylogeny (ParSNP/Roary) to
test whether Kalidasan's environmental isolates cluster in a known Sgn3/Sgn4
subclade with characterised gene losses.

\subsection*{Q4 --- Transferrin as top iron source without transferrin receptors: mechanism?}
\textbf{Basis.} The paper's Fig~3 shows transferrin gives maximal growth
($p<0.001$), yet the paper's own 17-target Table~2 (which this replication
re-confirmed) contains \emph{no} TbpA/TbpB/LbpA/LbpB homologs. The mechanism
by which \textit{S.\ maltophilia} liberates iron from Tf is unaddressed
in the paper's Discussion, which waves at "siderophore could potentially
scavenge iron-saturated Tf" but does not test it. Positive control
\textit{N.\ meningitidis} \emph{does} have TbpAB/LbpAB and is not comparable.

\textbf{Next steps.} (1) HMMER search of all 4 genomes for TbpA (PF01298),
TbpB (PF07298), LbpA/LbpB, and any TonB-dependent receptor with predicted
Tf-binding specificity. (2) Growth on holo-Tf $\pm$ apoferritin controls;
if only iron-saturated citrate supports growth, \textit{S.\ maltophilia} is
not using Tf directly. (3) Knock out the secreted metalloproteases (StmPr1/2/3)
in K279a and repeat Tf-growth assay --- proteolytic iron release hypothesis.
(4) Cross-reference Garcia et al.\ 2015 (\textit{Res Microbiol}) on
\textit{S.\ maltophilia} lack of Tf receptors.

\subsection*{Q5 --- Are any of the 6 iron-induced targets required for virulence?}
\textbf{Basis.} The paper's own concluding sentence: "\textit{These data need
to be further confirmed through several knockout studies}" --- the authors
themselves flag it. As of 2026, no follow-up study has closed the loop for
the 6-target set (FeSR, HmuT, Hup, ETFb, TonB, Fur). This replication cannot
close it in silico; but it does re-confirm the genes exist in K279a and are
ready for mutagenesis. Also: the ETFb induction (2.28$\times$) is unexplained
mechanistically --- ETFb is an electron-transfer flavoprotein, not an
obviously iron-acquisition gene.

\textbf{Next steps.} (1) Markerless in-frame deletion mutants in K279a for
each target (pEX18Tc adaptation, sacB counter-selection). (2) Phenotype panel:
BHI-DIP growth, CAS/Arnow's siderophore production, biofilm formation, serum
survival (50\% human serum), \textit{Galleria mellonella} or mouse-lung
colonisation. (3) In-trans complementation to confirm phenotype specificity.
(4) For ETFb specifically: dual transcriptomics + metabolomics under iron
depletion in wildtype vs $\Delta$etfB --- test electron-transfer-vs-iron
hypothesis.
